\section{Dualitas Ruang Banach}

\begin{conv}\label{conv_subsp_Bsp} Dalam konteks ruang Banach, kita menggunakan konvensi yang sama untuk pemakaian kata
``subruang'' seperti yang kita gunakan untuk ruang Hilbert. Kata ``subruang'' akan selalu berarti
 \index{subruang!dari ruang Banach}%
 \index{konvensi!subruang ruang Banach bersifat tertutup}%
\emph{subruang vektor tertutup}. Untuk menunjukkan bahwa $M$ adalah subruang dari suatu ruang Banach $B$, kita menulis
 \index{<binrelle@$\preccurlyeq$ (subruang dari ruang Banach atau Hilbert)}%
$M \preccurlyeq B$.
\end{conv}

\begin{prop} Jika $M$ adalah subruang dari suatu ruang Banach $B$,
 \index{hasil bagi!ruang Banach}%
 \index{Banach!ruang!hasil bagi}%
 \index{ruang!hasil bagi!Banach}%
maka ruang linear bernorma hasil bagi $B/M$ juga merupakan ruang Banach.
\end{prop}

\begin{proof}[\emph{Petunjuk pembuktian}] Untuk menunjukkan bahwa $B/M$ lengkap, cukup ditunjukkan bahwa dari
sebarang barisan Cauchy $(u_n)$ dalam $B/M$ kita dapat mengambil suatu subbarisan yang konvergen. Untuk itu,
perhatikan bahwa jika $(u_n)$ adalah barisan Cauchy, maka untuk setiap $k \in \N$ terdapat $N_k \in \N$ sedemikian sehingga
$\norm{u_n - u_m} < 2^{-k}$ setiap kali $m$, $n \ge N_k$. Sekarang gunakan induksi. Pilih $n_1 = N_1$.
Setelah $n_1$, $n_2$, \dots, $n_k$ dipilih sedemikian sehingga $n_1 < n_2 < \dots < n_k$ dan $N_j \le n_j$ untuk
$1 \le j \le k$, pilih $n_{k+1} = \max\{N_{k+1}, n_k + 1\}$. Dengan demikian diperoleh subbarisan
$\bigl(u_{n_k}\bigr)$ yang memenuhi $\norm{u_{n_k} - u_{n_{k+1}}} < 2^{-k}$ untuk setiap $k \in \N$.
Dari setiap kelas ekuivalensi $u_{n_k}$, pilih suatu vektor $x_k$ dalam $B$ sedemikian sehingga $\norm{x_k - x_{k+1}} < 2^{-k+1}$.
Kemudian Proposisi~\ref{prop_abs_sum_implies_sum} menunjukkan bahwa barisan $(x_k - x_{k+1})$ dapat
dijumlahkan. Dari sini, simpulkan bahwa barisan $(x_k)$ dan $\bigl(u_{n_k}\bigr)$ konvergen.
(Rincian argumen ini dituliskan dengan baik dalam pembuktian Teorema 3.4.13 pada~\cite{BrownP:1970}.) \ns
\end{proof}

\begin{prop} Jika $J$ adalah ideal tertutup dalam suatu aljabar Banach $A$,
 \index{hasil bagi!aljabar Banach}%
 \index{Banach!aljabar!hasil bagi}%
 \index{aljabar!hasil bagi!Banach}%
maka $A/J$ merupakan aljabar Banach.
\end{prop}

\begin{thm}[Teorema hasil bagi fundamental untuk $\cat{BALG}$] Misalkan $A$ dan $B$ aljabar Banach
 \index{BALG@$\cat{BALG}$!hasil bagi dalam}%
 \index{hasil bagi!teorema!untuk $\cat{BALG}$}%
dan $J$ ideal tertutup sejati dalam~$A$. Jika $\phi$ suatu homomorfisme dari $A$ ke $B$ dan
$\ker\phi \supseteq J$, maka terdapat homomorfisme tunggal $\wt\phi\colon A/J \sto
B$ yang membuat diagram berikut komutatif.
   \[ \xymatrix@+30pt{A \ar[d]_*+{\pi} \ar[dr]^*+{\phi} & \\
                           A/J \ar[r]_*+{\widetilde\phi} & B  }\]
Selain itu, $\wt\phi$ injektif jika dan hanya jika $\ker\phi = J$; dan $\wt\phi$
surjektif jika dan hanya jika $\phi$ surjektif.
\end{thm}

\begin{prop}\label{C037821} Misalkan $J$ ideal tertutup sejati dalam suatu aljabar Banach komutatif
beridentitas~$A$. Maka $J$ maksimal jika dan hanya jika $A/J$ adalah medan.
\end{prop}

Proposisi berikut menyatakan syarat-syarat yang cukup sederhana agar suatu epimorfisme antara ruang Banach
ternyata menjadi faktor dalam morfisme-morfisme lain. Sepintas lalu, kegunaan hasil semacam ini
mungkin tampak agak meragukan. Namun, ketika membuktikan Proposisi~\ref{00151231}, perhatikan bahwa
inilah \emph{tepatnya} yang Anda perlukan.

\begin{prop}\label{prop_surj_as_factor} Tinjau kategori ruang Banach
dan pemetaan linear kontinu. Jika $T\colon A \sto C$, $R\colon A \sto B$, $R$
surjektif, dan $\ker R \subseteq \ker T$, maka terdapat
$S\colon B \sto  C$ tunggal sedemikian sehingga $SR = T$. Selain itu, $S$ injektif jika dan
hanya jika $\ker R = \ker T$.
\end{prop}

\begin{proof}[\emph{Petunjuk pembuktian}] Tinjau diagram komutatif berikut.
Gunakan \emph{teorema pemetaan terbuka}.
 \[\xy
    \scalefactor{1.4}%
    \Atrianglepair/>`>`>`<-`>/[A`B`A/\ker R`C;R`\pi`T`\wt R`\wt T]
 \endxy\]   \ns
\end{proof}

Menarik untuk diperhatikan bahwa hasil sebelumnya dapat dinyatakan dengan cara yang
sedikit memperkuat klaim keunikannya. Sebagaimana dinyatakan, kesimpulannya ialah bahwa terdapat
pemetaan linear terbatas $S$ tunggal sedemikian sehingga $SR = T$. Tanpa kerja tambahan, kita dapat merumuskan ulang
proposisi tersebut sebagai berikut.

\begin{prop}\label{prop_surj_as_factor2} Tinjau kategori ruang Banach
dan pemetaan linear kontinu. Jika $T\colon A \sto C$, $R\colon A \sto B$, $R$
surjektif, dan $\ker R \subseteq \ker T$, maka terdapat fungsi tunggal
$S$ yang memetakan $B$ ke $C$ sedemikian sehingga $S\circ R = T$. Fungsi $S$ harus
terbatas dan linear. Selain itu, fungsi tersebut injektif jika dan hanya jika $\ker R = \ker T$.
\end{prop}

Apakah versi kedua hasil ini, yang lebih terperinci, layak dinyatakan? Menurut saya, jawabannya bergantung pada
arah yang hendak dituju. Memilih versi yang lebih panjang hanya dengan alasan bahwa versi itu lebih umum
tampaknya tidak menarik. Di pihak lain, andaikan kemudian kita memerlukan hasil berikut.

\begin{prop}\label{prop_factor_blm} Jika suatu pemetaan linear terbatas $T\colon A \sto C$ antara ruang-ruang Banach dapat difaktorkan
menjadi komposisi dua fungsi $T = f \circ R$, dengan $R \in \ofml B(A,B)$ surjektif,
maka $f \colon B \sto C$ terbatas dan linear.
\end{prop}

\begin{proof}[\emph{Petunjuk pembuktian}] Untuk menerapkan~\ref{prop_surj_as_factor2}, yang diperlukan di sini hanyalah
memperhatikan bahwa $f(0) = 0$ sehingga $\ker R \subseteq \ker T$.   \ns
\end{proof}

Karena hasil sebelumnya merupakan akibat yang hampir langsung dari~\ref{prop_surj_as_factor2}, ada baiknya
versi yang lebih terperinci ini tersedia. Orang mungkin berpendapat bahwa meskipun~\ref{prop_factor_blm} barangkali tidak
sepenuhnya jelas dari pernyataan~\ref{prop_surj_as_factor}, hasil itu tentu jelas dari
pembuktiannya, jadi untuk apa memakai versi yang lebih rumit? Hal-hal semacam ini tentu merupakan soal selera.

\begin{defn}\label{defn_Bsp_exact_seq} Suatu barisan ruang Banach dan pemetaan linear kontinu
   \[\xymatrix{{\cdots}\ar[r] & B_{n-1}\ar[r]^{T_n}
          & B_n\ar[r]^{T_{n+1}} & B_{n+1}\ar[r] & {\cdots}
   }\]
dikatakan
 \index{barisan!eksak}%
 \index{eksak!barisan}%
\df{eksak di} $B_n$ jika $\ran T_n = \ker T_{n+1}$. Suatu barisan disebut \df{eksak} jika barisan tersebut
eksak di setiap ruang Banach penyusunnya. Suatu barisan ruang Banach dan
pemetaan linear kontinu berbentuk
   \begin{equation}\label{001500202i}\xymatrix{
      0 \ar[r] & A \ar[r]^S & B\ar[r]^T & C\ar[r] & 0
   }\end{equation}
 \index{barisan!eksak pendek}%
 \index{barisan eksak pendek}%
yang eksak di $A$, $B$, dan $C$ disebut \df{barisan eksak pendek}. (Di sini $\vc 0$ menyatakan ruang Banach trivial
berdimensi nol, dan panah-panah tanpa label adalah pemetaan yang jelas.)
\end{defn}

Definisi-definisi sebelumnya diberikan untuk
 \index{BAN@$\cat{BAN_\infty}$!kategori}%
 \index{kategori!$\cat{BAN_\infty}$ sebagai suatu}%
kategori $\cat{BAN_\infty}$ dari ruang Banach dan pemetaan linear kontinu. Tentu saja,
gagasan tentang suatu \emph{barisan eksak} bermakna dalam banyak situasi---misalnya, dalam
kategori aljabar Banach, aljabar-$C^*$, ruang Hilbert, ruang vektor, grup Abel,
dan modul, di antara yang lainnya.

\begin{prop}\label{00151231} Tinjau diagram berikut dalam kategori ruang Banach
dan pemetaan linear kontinu
 \[\xymatrix{
      0\ar[r] & A\ar[d]^f\ar[r]^j & B\ar[d]^g\ar[r]^k
              & C\ar@{-->}[d]^h \ar[r] & 0 \\
      0\ar[r] & A'\ar[r]_{j'} & B'\ar[r]_{k'}
      & C'\ar[r] & 0
 }\]
Jika baris-barisnya eksak dan persegi kirinya komutatif, maka terdapat pemetaan linear kontinu
$h\colon C \sto C'$ tunggal yang membuat persegi kanannya komutatif.
\end{prop}

\begin{prop}[Lema Lima]\label{00151232}
 \index{lema lima}%
Andaikan dalam diagram ruang Banach dan pemetaan linear kontinu berikut
 \[\xymatrix{
      0\ar[r] & A\ar[d]^f\ar[r]^j & B\ar[d]^g\ar[r]^k
              & C\ar[d]^h \ar[r] & 0 \\
      0\ar[r] & A'\ar[r]_{j'} & B'\ar[r]_{k'}
      & C'\ar[r] & 0
 }\]
baris-barisnya eksak dan persegi-perseginya komutatif. Maka pernyataan-pernyataan berikut berlaku.
 \begin{enumerate}[label=\rm{(\alph*)}]
  \item Jika $g$ surjektif, maka $h$ juga surjektif.
  \item Jika $f$ surjektif dan $\ran j \supseteq \ker g$, maka $h$ injektif.
  \item Jika $f$ dan $h$ surjektif, maka $g$ juga surjektif.
  \item Jika $f$ dan $h$ injektif, maka $g$ juga injektif.
 \end{enumerate}
\end{prop}

Ada banyak keuntungan bekerja dengan operator pada ruang Hilbert atau Banach yang memiliki jangkauan tertutup.
Salah satu keuntungan langsung ialah bahwa kita memperoleh kembali kesederhanaan \emph{teorema dasar aljabar linear}
(\ref{00016025}) dari versi-versi yang agak lebih rumit dalam Proposisi~\ref{prop_ker_adj_perp_ran_op}
dan Teorema~\ref{prop_FTLA_Bsp}. Di pihak lain, ada pula kerugiannya. Operator-operator itu tidak membentuk suatu kategori; komposisi
dua operator berjangkauan tertutup dapat gagal memiliki jangkauan tertutup.

\begin{prop} Jika suatu pemetaan linear terbatas $T\colon B \sto C$ antara dua ruang Banach memiliki jangkauan tertutup, maka
adjoinnya~$T^*$ juga demikian dan $\ran T^* = (\ker T)^\perp$.
\end{prop}

\begin{proof}[\emph{Petunjuk pembuktian}] Simpulkan dari Proposisi~\ref{prop_FTLA_Bsp}(d) bahwa kita hanya perlu menunjukkan bahwa
$(\ker T)^\perp \subseteq \ran T^*$. Misalkan $f$ suatu fungsional sebarang dalam $(\ker T)^\perp$. Yang harus dicari adalah suatu
fungsional $g$ dalam $C^*$ sedemikian sehingga $f = T^*g$. Untuk itu, gunakan \emph{teorema pemetaan terbuka} untuk menunjukkan bahwa pemetaan $\wt T_\ast$
dalam diagram berikut invertibel. (Pemetaan $T_\ast\colon B \sto \ran T$ dan $T$ sama persis, kecuali bahwa
kodomainnya berbeda.)

 \[\xy
    \scalefactor{1.4}%
    \Atrianglepair/>`>`>`<-`>/[B`\ran T`B/\ker T`\K;T_\ast`\pi`f`\wt{T_\ast}`\wt f]
 \endxy\]

Misalkan $g_0 := \wt f\,{\wt T_\ast}^{-1}$. Gunakan \emph{teorema Hahn--Banach} untuk menghasilkan~$g$ yang diinginkan.   \ns
\end{proof}

\begin{prop} Dalam kategori $\cat{BAN_\infty}$, jika barisan $\xymatrix{A \ar[r]^S & B \ar[r]^T & C}$ eksak di $B$
dan $T$ memiliki jangkauan tertutup, maka $\xymatrix{C^* \ar[r]^{T^*} & B^* \ar[r]^{S^*} & A^*}$ eksak di~$B^*$.
\end{prop}

\begin{defn} Suatu funktor dari kategori ruang Banach dan pemetaan linear kontinu ke dirinya sendiri disebut
 \index{eksak!funktor}%
 \index{funktor!eksak}%
\df{eksak} jika funktor tersebut membawa barisan eksak pendek ke barisan eksak pendek.
\end{defn}

\begin{exam}\label{exam_dualfunctor_exact} Funktor $B \mapsto B^*$, $T \mapsto T^*$ (lihat~\ref{exam_dual_functor})
dari kategori ruang Banach dan pemetaan linear kontinu ke dirinya sendiri bersifat eksak.
\end{exam}

\begin{exam} Misalkan $B$ suatu ruang Banach dan $\sbsb{\tau}B\colon x \mapsto x^{**}$ pembenaman alami dari $B$ ke~$B^{**}$
(lihat Definisi~\ref{defn_nat_embed}). Maka ruang hasil bagi $B^{**}/\ran\sbsb{\tau}B$ merupakan ruang Banach.
\end{exam}

\begin{notn} Kita akan menyatakan dengan $\clo B$ ruang hasil bagi
$B^{**}/\ran \sbsb{\tau}B$ dalam contoh sebelumnya.
\end{notn}

\begin{prop} Jika $T\colon B \sto C$ suatu pemetaan linear kontinu antara ruang-ruang Banach, maka terdapat pemetaan linear kontinu tunggal
$\clo T\colon \clo B \sto \clo C$ sedemikian sehingga $\clo T(\phi + \ran\sbsb{\tau}B) = (T^{**}\phi) + \ran\sbsb{\tau}C$ untuk setiap $\phi \in B^{**}$.
\end{prop}

\begin{proof}[\emph{Petunjuk pembuktian}] Terapkan~\ref{00151231} pada diagram berikut:
 \begin{equation}\label{eqn_exactCD_Bbar}\xymatrix{
      0\ar[r] & B\ar[d]^T\ar[r]^{\sbsb{\tau}B} & B^{**}\ar[d]^{T^{**}}\ar[r]^{\sbsb{\pi}B}
              & \clo B\ar@{-->}[d]^{\clo T} \ar[r] & 0 \\
      0\ar[r] & C\ar[r]_{\sbsb{\tau}C} & C^{**}\ar[r]_{\sbsb{\pi}C} & \clo C\ar[r] & 0
 }\end{equation}  \ns
\end{proof}

\begin{exam}\label{exam_bar_functor} Pasangan pemetaan $B \mapsto \clo B$, $T \mapsto \clo T$ dari
kategori $\cat{BAN_\infty}$ ke dirinya sendiri merupakan suatu
 \index{<unopu@$B \mapsto \clo B$, $T \mapsto \clo T$ (funktor)}%
 \index{funktor!dari suatu ruang Banach $B$ ke $\clo B$}%
funktor kovarian eksak.
\end{exam}

\begin{prop} Dalam kategori ruang Banach dan pemetaan linear kontinu, jika
barisan
 \[\begin{CD}
  0 \longrightarrow A @> S >> B @> T >> C \longrightarrow 0
  \end{CD}\]
eksak, maka diagram berikut komutatif dan memiliki baris-baris
serta kolom-kolom yang eksak.
 \[\bfig
    \iiixiii{'7777}[A`A^{**}`\overline A`B`B^{**}`\overline B`C`C^{**}`\overline C;%
``````S`{S^{**}}`{\overline S}`T`{T^{**}}`{\overline T}]
  \efig\]
\end{prop}

\begin{cor} Setiap subruang dari suatu ruang Banach refleksif bersifat refleksif.
\end{cor}

\begin{cor} Jika $M$ suatu subruang dari ruang Banach refleksif~$B$, maka $B/M$ refleksif.
\end{cor}

\begin{cor} Misalkan $M$ suatu subruang dari ruang Banach~$B$. Jika $M$ dan $B/M$ keduanya refleksif,
maka $B$ juga refleksif.
\end{cor}



















\begin{defn} Misalkan $T \in \ofml L(B,C)$, dengan $B$ dan $C$ ruang Banach. Definisikan
$\coker T$, yaitu
 \index{kokernel}%
 \index{kokernel@$\coker T$ (kokernel dari $T$)}%
\df{kokernel} dari~$T$, sebagai $C/\ran T$.
\end{defn}



Menentukan secara eksplisit dual dari suatu ruang Banach sering kali agak menantang. Berikut diberikan beberapa
contoh penting. Namun, pertama-tama perlu disampaikan satu peringatan. Notasi untuk objek-objek dalam ruang linear bernorma
yang terdiri atas barisan dapat menyesatkan. Ketika membahas kelengkapan ruang-ruang semacam itu, kita berhadapan dengan barisan dari
barisan bilangan kompleks. Demi memudahkan pembaca (dan mungkin juga diri Anda sendiri sebulan kemudian,
ketika Anda mencoba membaca sesuatu yang Anda tulis!), sebaiknya bedakan secara notasi antara barisan bilangan kompleks
dan barisan dari barisan. Ada banyak cara untuk melakukannya: subskrip ganda, subskrip dan
superskrip, notasi fungsional untuk barisan kompleks, \emph{dan sebagainya}.

Notasi bermasalah lain yang muncul dalam ruang barisan adalah $\sum_{k=1}^\infty a_k \vc e^k$, dengan $\vc e^k$
merupakan barisan yang bernilai $1$ pada koordinat ke-$k$ dan $0$ di semua koordinat lainnya. (Ingat bahwa dalam
Contoh~\ref{C063149}, kita menggunakan vektor-vektor ini sebagai basis ortonormal biasa (atau standar) untuk ruang Hilbert
$l_2$.) Tampaknya hampir tidak ada kebingungan yang timbul jika $\sum_{k=1}^\infty a_k \vc e^k$ dipandang
hanya sebagai singkatan untuk barisan bilangan kompleks $(a_k) = (a_1,a_2,a_3,\dots)$. Namun, jika seseorang hendak
memperlakukannya secara harfiah sebagai deret tak hingga, diperlukan kehati-hatian. Dalam ruang $\sbsb c0$ dan $l_1$, tidak ada
masalah. Deret itu memang merupakan limit, ketika $n$ membesar, dari $\sum_{k=1}^n a_k \vc e^k$. (Mengapa?) Namun, dalam ruang
$c$ dan $l_\infty$, misalnya, terjadi kegagalan yang sangat buruk. Untuk melihat apa yang dapat terjadi, tinjau barisan konstan
$(1,1,1,\dots)$ (yang merupakan anggota $c$ maupun $l_\infty$). Maka
$\lim_{n \sto\infty}\norm{(1,1,1,\dots) - \sum_{k=1}^n \vc e^k}$ gagal mendekati nol dalam kedua kasus tersebut.
Dalam $c$ maupun $l_\infty$, limitnya adalah~$1$.

\begin{exam}\label{exam_dual_C0} Ruang linear bernorma $\sbsb c0$ (lihat~\ref{C0231303})
 \index{c@$c_0$!sebagai ruang Banach}%
 \index{Banach!ruang!$c_0$ sebagai suatu}%
 \index{l@$l_1$!sebagai dual dari $c_0$}%
 \index{dual!dari $c_0$}%
merupakan ruang Banach dan dualnya adalah $l_1$ (lihat~\ref{C023191}).
\end{exam}

\begin{proof}[\emph{Petunjuk pembuktian}] Tinjau pemetaan
$T\colon {\sbsb c0}^* \sto l_1\colon f \mapsto \sum_{k=1}^\infty f(\vc e^k)\vc e^k$.
(Jangan lupa menunjukkan bahwa $T$ benar-benar memetakan ke~$l_1$.) Periksa bahwa $T$ merupakan isomorfisme isometrik.
Untuk membuktikan bahwa $T$ surjektif, mulailah dengan suatu vektor $b \in l_1$ dan tinjau fungsi
$f\colon \sbsb c0 \sto \C\colon a \mapsto \sum_{k=1}^\infty a_kb_k$.    \ns
\end{proof}

\begin{cor} Ruang $l_1$ merupakan suatu
 \index{l@$l_1$!sebagai ruang Banach}%
 \index{Banach!ruang!$l_1$ sebagai suatu}%
ruang Banach.
\end{cor}

Menarik bahwa ternyata ruang $\sbsb c0$ dan $c$ tidak isomorfik secara isometrik,
meskipun ruang-ruang dualnya isomorfik secara isometrik. Perhitungan dual dari $c$ secara teknis sedikit lebih rumit daripada
perhitungan~${\sbsb c0}^{\!*}$. Mungkin berguna untuk meninjau \emph{basis Schauder} bagi kedua ruang ini.

\begin{defn} Suatu barisan $(\vc e^k)$ dari vektor-vektor dalam suatu ruang Banach $B$ disebut
 \index{basis Schauder}%
 \index{basis!Schauder}%
\df{basis Schauder} bagi ruang tersebut jika untuk setiap $x \in B$ terdapat barisan skalar $(\alpha_k)$ yang unik
sedemikian sehingga $x = \sum_{k=1}^\infty \alpha_k \vc e^k$.
\end{defn}

Dua catatan mengenai definisi sebelumnya: Kata ``unik'' sangat penting; dan tidak benar bahwa
setiap ruang Banach memiliki basis Schauder.

\begin{exam} Suatu basis ortonormal dalam ruang Hilbert separabel merupakan basis Schauder bagi ruang tersebut. (Lihat
Proposisi~\ref{prop_unique_onbasis} dan~\ref{prop_Schauder_basis_Hsp}.)
\end{exam}

\begin{exam}\label{exam_dual_c0} Kita memilih
 \index{biasa!basis Schauder untuk~$\sbsb c0$}%
 \index{Schauder!basis!biasa (atau standar) untuk~$\sbsb c0$}%
 \index{basis!standar (atau biasa)}%
vektor-vektor yang disebut \df{standar} (atau
 \index{standar!basis Schauder untuk $\sbsb c0$}%
\df{biasa}) sebagai \df{vektor-vektor basis} untuk $\sbsb c0$ dengan cara yang sama seperti untuk~$l_2$ (lihat~\ref{C063149}).
Untuk setiap $n \in \N$, misalkan $\vc e^n$ adalah barisan dalam $l_2$ yang koordinat ke-$n$-nya bernilai $1$
dan semua koordinat lainnya bernilai~$0$. Maka $\{\vc e^n\colon n \in \N\}$ merupakan basis Schauder
bagi~$\sbsb c0$.
\end{exam}

Sepintas lalu, ruang $c$ dari semua barisan konvergen mungkin tampak jauh lebih besar daripada $\sbsb c0$,
yang hanya memuat barisan-barisan yang konvergen ke nol. Lagi pula, untuk setiap vektor
$x$ dalam $\sbsb c0$ terdapat tak terhitung banyaknya vektor berbeda dalam $c$ (cukup tambahkan sebarang bilangan kompleks
tak nol pada setiap suku barisan~$x$). Mungkinkah $c$, yang memuat tak terhitung banyaknya
salinan dari koleksi anggota~$\sbsb c0$ yang tak terhitung, memiliki basis Schauder (yang mesti terhitung)?
Setelah direnungkan, jawabannya adalah \emph{ya}. Kita hanya perlu menambahkan satu elemen
pada basis bagi~$\sbsb c0$ untuk menangani translasi.

\begin{exam} Misalkan $\vc e^1$, $\vc e^2$, \dots seperti dalam Contoh~\ref{exam_dual_c0} sebelumnya. Selain itu,
nyatakan dengan $\vc 1$ barisan konstan $(1,1,1,\dots)$ dalam~$c$. Maka $(\vc 1, \vc e^1, \vc e^2, \dots)$
merupakan basis Schauder bagi~$c$.
\end{exam}

\begin{exam}\label{exam_dual_c} Ruang linear bernorma $c$ (lihat~\ref{C0231302})
 \index{c@$c$!sebagai ruang Banach}%
 \index{Banach!ruang!$c$ sebagai suatu}%
 \index{l@$l_1$!sebagai dual dari $c$}%
 \index{dual!dari $c$}%
merupakan ruang Banach dan dualnya adalah $l_1$ (lihat~\ref{C023191}).
\end{exam}

\begin{proof}[\emph{Petunjuk pembuktian}]
\leavevmode\par\noindent
Pertama-tama, perhatikan bahwa untuk setiap $f \in c^*$, barisan
$\bigl(f(\vc e^k)\bigr)$ dapat dijumlahkan secara mutlak dan, sebagai akibatnya, $\sum_{k=1}^\infty a_k f(\vc e^k)$
konvergen untuk setiap barisan $a = (a_1,a_2,\dots)$ dalam~$c$. Perhatikan pula bahwa setiap $a \in c$ dapat dituliskan sebagai
$a = a_\infty \vc 1 + \sum_{k=1}^\infty (a_k -a_\infty)\vc e^k$, dengan $a_\infty := \lim_{k \sto \infty}a_k$
dan $\vc 1 := (1,1,1,\dots)$. Gunakan hal ini untuk membuktikan bahwa untuk setiap $f \in c^*$,
   \[ \abs{f(a)} \le \biggl[ \abs{\beta_f} + \sum_{k=1}^\infty \abs{f(\vc e^k)} \biggr] {\norm a}_\infty \]
dengan $\beta_f := f(\vc 1) - \sum_{k=1}^\infty f(\vc e^k)$. Untuk $t = (t_0,t_1,t_2,\dots) \in l_1$ dan
$a = (a_1,a_2,a_3,\dots) \in c$, definisikan
   \[ S_t(a) := t_0 a_\infty + \sum_{k=1}^\infty t_ka_k \]
dan periksa bahwa $S\colon t \mapsto S_t$ merupakan pemetaan linear kontinu dari $l_1$ ke~$c^*$.
Selanjutnya, untuk $f \in c^*$, definisikan
   \[ Tf := (\beta_f, f(\vc e^1), f(\vc e^2), \dots) \]
dan tunjukkan bahwa $T\colon f \mapsto Tf$ merupakan isometri linear dari $c^*$ ke~$l_1$. Sekarang, untuk melihat bahwa $c^*$ dan
$l_1$ isomorfik secara isometrik, cukup ditunjukkan bahwa $ST = \id{c^*}$ dan $TS = \id{l_1}$.  \ns
\end{proof}












\begin{defn} Suatu ruang topologis disebut
 \index{lokal!kompak}%
 \index{kompak!lokal}%
\df{kompak lokal} jika setiap titik dalam ruang tersebut memiliki lingkungan kompak.
\end{defn}

\begin{notn} Misalkan $X$ suatu ruang kompak lokal. Suatu fungsi bernilai skalar $f$ pada $X$ termasuk dalam $\fml C_0(X)$ jika
 \index{C@$\fml C_0(X)$!ruang fungsi kontinu yang lenyap di tak hingga}%
fungsi tersebut kontinu dan lenyap di tak hingga (yakni, jika untuk setiap bilangan $\epsilon > 0$ terdapat subhimpunan kompak
dari $X$ yang di luarnya berlaku $\abs{f(x)} < \epsilon$).
\end{notn}

\begin{exam}\label{exam_C0X} Jika $X$ suatu ruang Hausdorff kompak lokal, maka, terhadap operasi aljabar titik demi titik dan norma seragam,
 \index{Banach!ruang!$\fml C_0(X)$ sebagai suatu}%
 \index{C@$\fml C_0(X)$!sebagai ruang Banach}%
$\fml C_0(X)$ merupakan ruang Banach.
\end{exam}

Kita mengingat kembali salah satu teorema terpenting dari analisis real.

\begin{thm}[Representasi Riesz]\label{C057951} Misalkan $X$ suatu ruang Hausdorff kompak lokal.
 \index{Riesz!teorema representasi}%
Jika $\phi$ suatu fungsional linear terbatas pada $\fml C_0(X)$, maka terdapat ukuran Borel kompleks reguler
$\mu$ yang unik pada $X$ sedemikian sehingga
  \[ \phi(f) = \int_X f\,d\mu \]
untuk setiap $f \in \fml C_0(X)$ dan $\norm\mu = \norm\phi$.
\end{thm}

\begin{proof} Teorema ini dan akibat berikutnya disajikan dengan baik dalam Lampiran C pada buku analisis
fungsional karya Conway~\cite{Conway:1990}. Lihat juga~\cite{Folland:1984}, Teorema 7.17; \cite{HewittS:1965},
Teorema 20.47 dan 20.48; \cite{McDonaldW:1999}, Teorema 9.16; dan \cite{Rudin:1987}, Teorema 6.19.  \ns
\end{proof}

\begin{cor}\label{C057952} Jika $X$ suatu ruang Hausdorff kompak lokal, maka $\bigl(\fml C_0(X)\,\bigr)^*$
 \index{dual!dari $\fml C_0(X)$}%
 \index{C@$\fml C_0(X)$!dual dari}%
 \index{M@$\textbf{M}(X)$!sebagai dual dari $\fml C_0(X)$}%
isomorfik secara isometrik dengan~$\textbf{M}(X)$,
 \index{M@$\textbf{M}(X)$!ukuran Borel kompleks reguler pada~$X$}%
ruang Banach dari semua ukuran Borel kompleks reguler pada~$X$.
\end{cor}

\begin{cor} Jika $X$ suatu ruang Hausdorff kompak lokal, maka $\textbf{M}(X)$
 \index{M@$\textbf{M}(X)$!sebagai ruang Banach}%
 \index{Banach!ruang!$\textbf{M}(X)$ sebagai suatu}%
merupakan ruang Banach.
\end{cor}










Dalam~\ref{defn_nls_adjoint}, kita mendefinisikan adjoin dari suatu operator linear antara ruang-ruang Banach,
dan dalam~\ref{def000926}, kita mendefinisikan adjoin dari suatu operator ruang Hilbert.
Dalam kasus operator pada ruang Hilbert (yang juga merupakan ruang Banach),
wajar untuk menanyakan hubungan apa, jika ada, di antara kedua ``adjoin'' ini.
Keduanya tentu tidak dapat sama, karena yang pertama bekerja di antara ruang-ruang dual, sedangkan yang kedua
di antara ruang-ruang semula. Dalam proposisi berikut, kita menggunakan antiisomorfisme
$\psi$ yang didefinisikan dalam~\ref{000341} (lihat~\ref{0022057}) untuk menunjukkan bahwa kedua adjoin itu
``pada hakikatnya'' sama.

\begin{prop} Misalkan $T$ suatu operator pada ruang Hilbert $H$ dan $\psi$ antiisomorfisme yang
didefinisikan dalam~\ref{000341}. Jika kita (untuk sementara) menyatakan dual ruang Banach
dari $T$ dengan $T'\colon H^* \sto H^*$, maka $T' = \psi\, T^* \psi^{-1}$. Artinya,
diagram berikut komutatif.
   \[ \xy
          \Square/{->}`{<-}`{<-}`{->}/[H^*`H^*`H`H;T'`\psi`\psi`T^*]
      \endxy \]
\end{prop}

\begin{defn}\label{0022091} Misalkan $A$ suatu subhimpunan dari ruang Banach $B$. Maka
 \index{anihilator}%
\df{anihilator} dari $A$, yang dinotasikan
 \index{<perp@$A^\perp$ (anihilator dari suatu himpunan)}%
dengan~$A^\perp$, adalah $\{f \in B^*\colon f(a) = 0 \text{ untuk setiap $a \in A$}\}$.
\end{defn}

Konflik notasi antara \ref{0001511} dan \ref{0022091} dapat
menimbulkan kebingungan. Untuk melihat bahwa dalam ruang Hilbert, anihilator suatu subruang dan
komplemen ortogonal subruang itu ``pada hakikatnya'' merupakan hal yang sama, gunakan
antiisomorfisme isometrik $\psi$ antara $H$ dan $H^*$ yang dibahas dalam~\ref{0022057} untuk mengidentifikasi keduanya.

\begin{prop} Misalkan $M$ suatu subruang dari ruang Hilbert $H$. Jika kita (untuk sementara) menyatakan
anihilator dari $M$ dengan~$M^a$, maka $M^a = \psi^{\sto}\bigl(M^\perp\bigr)$.
\end{prop}
